% LegalLean -- Verified Legal Reasoning
% Springer LNCS format
\documentclass[runningheads]{llncs}

\usepackage[T1]{fontenc}
\usepackage[utf8]{inputenc}
\usepackage{graphicx}
\usepackage{amssymb}
\usepackage{listings}
\usepackage{xcolor}

% ---- Lean 4 listings language definition ----
\definecolor{lean-keyword}{rgb}{0.13,0.29,0.53}
\definecolor{lean-comment}{rgb}{0.38,0.62,0.38}
\definecolor{lean-string}{rgb}{0.58,0.00,0.00}
\definecolor{lean-type}{rgb}{0.40,0.00,0.60}
\definecolor{lean-bg}{rgb}{0.97,0.97,0.99}

\lstdefinelanguage{lean4}{
  morekeywords={
    theorem, def, lemma, structure, inductive, class, instance,
    where, deriving, extends, with, by, do, let, fun, match, if, then, else,
    return, have, show, exact, apply, simp, constructor, intro, cases, induction,
    rfl, trivial, contradiction, decide, norm_num, ring, linarith,
    and, or, not, true, false, Type, Prop, Sort,
    forall, exists, namespace, end, open, section, variable, axiom,
    noncomputable, private, protected, mutual, partial, unsafe,
    List, Option, Bool, Nat, Int, String, Unit, Empty
  },
  sensitive=true,
  morecomment=[l]{--},
  morecomment=[s]{/-}{-/},
  morestring=[b]",
  morestring=[b]',
}

% Literate replacements for Unicode characters used in Lean 4 code
\lstset{
  language=lean4,
  basicstyle=\ttfamily\small,
  keywordstyle=\color{lean-keyword}\bfseries,
  commentstyle=\color{lean-comment}\itshape,
  stringstyle=\color{lean-string},
  backgroundcolor=\color{lean-bg},
  frame=single,
  framesep=4pt,
  rulecolor=\color{gray!40},
  breaklines=true,
  breakatwhitespace=false,
  showstringspaces=false,
  tabsize=2,
  captionpos=b,
  xleftmargin=6pt,
  xrightmargin=6pt,
  literate=
    {∀}{$\forall$}{1}
    {∃}{$\exists$}{1}
    {∧}{$\wedge$}{1}
    {∨}{$\vee$}{1}
    {→}{$\rightarrow$}{1}
    {←}{$\leftarrow$}{1}
    {↔}{$\leftrightarrow$}{1}
    {¬}{$\neg$}{1}
    {≠}{$\ne$}{1}
    {≥}{$\ge$}{1}
    {≤}{$\le$}{1}
    {⟨}{$\langle$}{1}
    {⟩}{$\rangle$}{1}
    {∉}{$\notin$}{1}
    {∈}{$\in$}{1}
    {·}{$\cdot$}{1}
    {α}{$\alpha$}{1}
    {β}{$\beta$}{1}
    {⊗}{$\otimes$}{1}
    {⇒}{$\Rightarrow$}{1}
    {⊢}{$\vdash$}{1}
    {ℕ}{$\mathbb{N}$}{1},
}

% Plain listing style for non-Lean code snippets (DDL etc.)
\lstdefinestyle{plain}{
  language={},
  basicstyle=\ttfamily\small,
  backgroundcolor=\color{lean-bg},
  frame=single,
  framesep=4pt,
  rulecolor=\color{gray!40},
  breaklines=true,
  showstringspaces=false,
  xleftmargin=6pt,
  xrightmargin=6pt,
  literate=
    {⊗}{$\otimes$}{1}
    {⇒}{$\Rightarrow$}{1}
    {∀}{$\forall$}{1}
    {∃}{$\exists$}{1}
    {≥}{$\ge$}{1}
    {≤}{$\le$}{1},
}

\begin{document}

\title{Verified Legal Reasoning: From Natural Language to Formal Proof}

\titlerunning{Verified Legal Reasoning}

\author{Eduardo Aguilar Pelaez\inst{1,2}}

\authorrunning{E. Aguilar Pelaez}

\institute{
  Legal Engine Ltd, London, UK\\
  \email{edu@legalengine.co.uk}
  \and
  Imperial College London, London, UK\\
  \email{e.aguilar@imperial.ac.uk}
}

\maketitle

\begin{abstract}
Legal AI systems increasingly serve as transducers -- converting unstructured
statutory text into structured representations -- but they lack a verification
layer capable of distinguishing what the type checker can certify from what
requires human judgement. This paper presents LegalLean, a framework for
verified legal reasoning built on Lean~4 dependent types. The central
contribution is the \texttt{FormalisationBoundary} type and the
\texttt{RequiresHumanDetermination} typeclass, which together encode Hart's
``open texture'' distinction as a first-class object in the type system rather
than an external limitation on formalisation. We demonstrate the framework
across three case studies: IRC~\S163(h) mortgage interest deduction, the O-1A
extraordinary ability visa (8~CFR~\S204.5(h)), and the Australian
Telecommunications Consumer Protections Code (C628:2012). The system produces
44 machine-verified theorems with zero unresolved \texttt{sorry} placeholders.
The headline result is \texttt{monotonicity\_break\_general}: a machine-checked
proof that the Tax Cuts and Jobs Act 2017 removed a previously valid deduction
permission for all home equity interest payments -- a formally verified instance
of non-monotonic statutory change. We compare LegalLean against four prior
systems (DDL, Catala, L4L, LegalRuleML) and discuss honest limitations.

\keywords{Dependent types \and Legal formalisation \and Lean~4 \and Open
texture \and Formal verification \and Defeasible reasoning}
\end{abstract}

% ---------------------------------------------------------------
\section{Introduction}
% ---------------------------------------------------------------

Google DeepMind's description of its language model as a ``transducer and
verifier'' articulates a productive division of labour: the language model
converts messy natural-language input into structured form, while a separate
verification layer checks the logic. This framing, applied to law, identifies
exactly what is missing from current legal AI systems. Systems that extract
rules from statutes, answer legal queries, or assess compliance are, in this
terminology, transducers. They produce structured representations. None of them
includes a verifier in the formal sense: a component that provides
machine-checkable guarantees about the reasoning it performs.

The gap matters because legal reasoning has a well-understood failure mode.
Hart's analysis in \textit{The Concept of Law}~\cite{hart1961concept} identified
that legal concepts have a core of settled meaning and a penumbra of
uncertainty. A formalisation that does not distinguish these zones does not
merely lack nuance -- it silently misrepresents where its own confidence ends.
When a compliance system returns ``non-compliant'', the user cannot tell whether
that verdict rests on a formally verifiable ground (the payment was not made by
an individual) or a contestable judgement (the property may not qualify as a
``qualified residence''). Both appear as outputs; only one is machine-checkable.

Existing approaches address parts of this problem but not the whole. Catala~\cite{merigoux2021catala}
is a functional programming language designed for statutes: it computes legal
outcomes deterministically but has no mechanism for representing open texture --
the code either handles a case or raises a runtime error. Governatori's
Defeasible Deontic Logic (DDL)~\cite{governatori2005representing,governatori2019compliance}
encodes legal rules as defeasible clauses with priority relations, handling
non-monotonicity well but providing no machine-checkable proofs and treating
vague predicates as ordinary literals. L4L (Logic for Law)~\cite{hou2021l4l}
uses dependent types in a style close to ours and produces SMT-backed proofs,
but the SMT layer is not fully transparent: the proof obligations are discharged
by a solver whose output is not an auditable term. LegalRuleML~\cite{palmirani2018modelling}
is a rich XML vocabulary for regulatory knowledge but provides no automated
reasoning whatsoever. None of these systems makes the formalisation boundary
itself a typed, auditable object.

LegalLean takes a different approach. It uses Lean~4~\cite{demoura2021lean4}
dependent types, where the Curry-Howard
correspondence~\cite{curry1934functionality,howard1980formulae} establishes that
propositions are types and proofs are terms. A legal conclusion -- ``this
interest payment is not deductible'' -- is a type. A proof of that conclusion is
a term of that type. If the term type-checks, the conclusion is verified; if it
does not, the error is explicit and localised. Crucially, when a legal condition
requires human determination (because the underlying predicate is vague or
discretionary), LegalLean does not attempt to discharge it. Instead, it wraps
the proposition in a \texttt{FormalisationBoundary} value that carries
machine-readable metadata about which institution must determine it and why.

An important architectural distinction: LegalLean is a verification layer, not an inference engine. The case study functions (such as \texttt{isDeductible} for IRC~\S163(h)) are formalisations of statutory logic written by a human or language model. They are programs to be verified, not derivations from a rule base. The value of Lean~4 is not that it replaces the formaliser but that it catches errors in the formalisation: if the encoded logic is inconsistent with the declared types, the type checker rejects it. A generic defeasible reasoning solver (in the style of DDL's proof theory) would be a valuable complement but serves a different role. DDL computes which conclusions follow from a set of rules; LegalLean verifies that a specific encoding of those rules is internally consistent and correctly typed. The two are complementary, not competitive.

The contributions of this paper are:
\begin{enumerate}
  \item The \texttt{FormalisationBoundary} inductive type and the
    \texttt{RequiresHumanDetermination} typeclass (with \texttt{Vague} and
    \texttt{Discretionary} subclasses), providing the first formal treatment of
    Hart's open texture as a machine-checkable type-level distinction.
  \item A verified encoding of IRC~\S163(h) including
    \texttt{monotonicity\_break\_general}, a machine-checked proof that the TCJA
    2017 constitutes a non-monotonic statutory change: for every home equity
    payment satisfying the pre-TCJA conditions, the legal outcome transitions
    from \texttt{boundary} (human determination required) to
    \texttt{formal~False} (machine-verified rejection).
  \item A dependent-type encoding of the O-1A visa ``3 of 8'' threshold using
    \texttt{AllDistinct} lists with length constraints, together with a
    \texttt{mixed} formalisation of the Kazarian two-part test.
  \item A structured comparison of LegalLean against
    DDL~\cite{governatori2005representing}, Catala~\cite{merigoux2021catala},
    L4L~\cite{hou2021l4l}, and LegalRuleML~\cite{palmirani2018modelling} across
    seven dimensions.
\end{enumerate}

All source code is available and compiles against Lean~4 v4.24.0 (the Aristotle
release) with no Mathlib dependency.

% ---------------------------------------------------------------
\section{Background}
% ---------------------------------------------------------------

\subsection{Hart's Open Texture}

In \textit{The Concept of Law} (1961), H.L.A. Hart~\cite{hart1961concept}
observed that legal concepts have a ``core of settled meaning'' and a
``penumbra of uncertainty''. The canonical example is a rule prohibiting
vehicles in a park: a car is unambiguously within the rule's core; a toy pedal
car is in the penumbra. The penumbral question cannot be resolved by logic alone
-- it requires a judgement about the rule's purpose and the case's particulars.
Hart called this irreducible property ``open texture'': the meaning of legal
terms cannot be fully determined in advance, and no formalisation can eliminate
the need for human judgement at the boundary.

Prior legal knowledge representation systems treat open texture as a limitation
imposed from outside: the formalisation goes as far as it can, and the rest is
left to the practitioner. LegalLean treats it as a property to be represented
inside the system.

\subsection{Default Logic and Defeasibility}

Legal reasoning is non-monotonic: adding new rules or evidence can retract
previously valid conclusions. Lawsky's ``A Logic for Statutes''~\cite{lawsky2017logic}
provides a careful analysis of how Reiter's default logic~\cite{reiter1980logic}
applies to statutory interpretation, demonstrating that statutes encode default
rules of the form: ``if conditions hold and no exception applies, the conclusion
follows.'' The exceptions are themselves defeasible, creating chains of defeat.

LegalLean represents defeasibility through the \texttt{Defeats} structure and
\texttt{DefeasiblyHolds} proposition type. Priority ordering (\textit{lex
posterior}, \textit{lex specialis}) is encoded in the \texttt{Priority} field of
each \texttt{LegalRule}. The key departure from prior work is that defeat
relations are typed: the \texttt{Defeats} structure carries a \texttt{reason}
field distinguishing \textit{lex posterior} from \textit{lex specialis}, and its
\texttt{strict} field distinguishes hard from soft defeat.

\subsection{Curry-Howard and Lean 4}

The Curry-Howard correspondence~\cite{curry1934functionality,howard1980formulae}
establishes an isomorphism between propositions and types: a proposition
\texttt{P} is inhabited (has a term of that type) if and only if it is provable.
In Lean~4~\cite{demoura2021lean4}, this means that constructing a value of type
\texttt{EvidenceFor P} is literally the same operation as proving \texttt{P}.
Lean~4's kernel verifier checks all proof terms, making the guarantee
machine-checkable in a stronger sense than SMT: the checker is a small,
auditable kernel, not an opaque solver.

Lean~4's dependent type system allows the type of a term to depend on its value.
This is essential for legal formalisation: the type of evidence required for
``meets 3 of 8 criteria'' depends on which specific criteria are claimed, and
\texttt{MeetsThreeCriteria} encodes this directly via an \texttt{AllDistinct}
list of \texttt{Criterion} values with a \texttt{cardinal} proof that
\texttt{satisfied.length >= 3}.

% ---------------------------------------------------------------
\section{The LegalLean Framework}
% ---------------------------------------------------------------

\subsection{Core Types}

The framework begins with three foundational types defined in
\texttt{LegalLean.Core}:

\begin{lstlisting}[language=lean4]
inductive Deontic where
  | permitted    : Deontic
  | prohibited   : Deontic
  | obligatory   : Deontic
  deriving Repr, BEq, Inhabited, DecidableEq
\end{lstlisting}

The deontic modalities cover the three fundamental normative positions. Lean~4's
\texttt{DecidableEq} derivation means that equality of modalities is decidable,
enabling automated conflict detection in theorems (for instance,
\texttt{R1\_BaseDeduction.modality~\textbackslash{}ne~R2\_TCJAProhibition.modality} is proved by
\texttt{simp}).

\begin{lstlisting}[language=lean4]
structure LegalRule (α : Type) where
  id          : String
  description : String
  modality    : Deontic
  priority    : Priority
  conclusion  : α → Prop
  defeatedBy  : List String
\end{lstlisting}

\texttt{LegalRule} is parameterised by the type \texttt{\(\alpha\)} of the
entity it governs: \texttt{LegalRule InterestPayment} for IRC~\S163(h) rules,
\texttt{LegalRule Provider} for TCP Code rules. The \texttt{conclusion} field is
a function from \texttt{\(\alpha\)} to \texttt{Prop}, encoding the rule's
conclusion as a dependent proposition. This is more expressive than DDL's flat
first-order signature.

\begin{lstlisting}[language=lean4]
structure EvidenceFor (P : Prop) where
  source   : String
  artifact : Option String
  witness  : P
\end{lstlisting}

\texttt{EvidenceFor P} is the propositions-as-types materialisation of legal
evidence: the \texttt{witness} field literally contains the proof of \texttt{P}.
Constructing an \texttt{EvidenceFor P} value is constructing a proof. The
\texttt{source} and \texttt{artifact} fields provide provenance that survives
code evolution -- they are not erased at compile time.

\subsection{The FormalisationBoundary Type}

\texttt{FormalisationBoundary} is the novel contribution of this framework:

\begin{lstlisting}[language=lean4]
inductive FormalisationBoundary (P : Sort _) where
  | formal   (proof : P)                                  : FormalisationBoundary P
  | boundary (reason : String) (authority : String)       : FormalisationBoundary P
  | mixed    (formalPart : String) (humanPart : String)   : FormalisationBoundary P
\end{lstlisting}

The three constructors correspond to three epistemic positions about a legal
proposition \texttt{P}:

\begin{itemize}
  \item \texttt{formal proof}: the type checker alone can verify \texttt{P}. No
    human determination is required. The \texttt{proof} term is a
    machine-checked certificate.
  \item \texttt{boundary reason authority}: \texttt{P} requires human
    determination. The \texttt{reason} string names the specific open-texture
    condition; \texttt{authority} names the institution responsible.
  \item \texttt{mixed formalPart humanPart}: \texttt{P} is partially formal and
    partially boundary. Part~1 of the Kazarian two-part test is \texttt{formal}
    (Lean verified the criteria count); Part~2 is \texttt{boundary} (USCIS
    holistic determination).
\end{itemize}

\subsection{The RequiresHumanDetermination Typeclass}

The \texttt{FormalisationBoundary} type marks where the boundary falls. The
\texttt{RequiresHumanDetermination} typeclass marks which propositions are
intrinsically at the boundary:

\begin{lstlisting}[language=lean4]
class RequiresHumanDetermination (P : Prop) where
  reason    : String
  authority : String

class Vague (P : Prop) extends RequiresHumanDetermination P where
  core      : String
  penumbra  : String

class Discretionary (P : Prop) extends RequiresHumanDetermination P where
  scope       : String
  constraints : List String
\end{lstlisting}

\texttt{Vague} captures Hart's core/penumbra distinction directly. \texttt{core}
describes settled cases; \texttt{penumbra} describes the contested region.
\texttt{Discretionary} captures a different phenomenon: not conceptual vagueness
but institutional authority. A USCIS adjudicator's determination of
``extraordinary ability'' is not vague in Hart's sense -- the adjudicator has
the power to decide. The typeclass distinguishes these two sources of
non-formalisability.

Both typeclasses constrain proof search: any proposition with a \texttt{Vague}
or \texttt{Discretionary} instance cannot be discharged by the type checker
without an explicit axiom. In the case studies, open-texture propositions such
as \texttt{HasExtraordinaryAbility} and \texttt{IsQualifiedResidence} are
introduced as axioms -- honest markers that the type checker cannot verify them.
The typeclass instances then carry the auditable metadata explaining why.

The use of axioms for open-texture propositions requires justification. An axiom
in Lean~4 is a proposition asserted without proof; it is not verified by the
kernel. One might object that this introduces a soundness gap: if the axioms are
inconsistent, any theorem built on them is vacuously true. This objection is
correct in principle but misdirected in practice. The axioms are not intended to
be discharged; they represent propositions that \textit{should not} be
machine-verifiable. A system that could prove \texttt{HasExtraordinaryAbility~a}
for a given applicant \texttt{a} would be claiming to have automated a
determination that the US Congress assigned to USCIS adjudicators. The axiom
marks the boundary between what the type checker should certify and what it
should defer. The theorems built on these axioms do not depend on their truth --
they depend on their \textit{structure}. The monotonicity break theorem, for
instance, does not use any open-texture axiom; it reasons entirely about the
formalisable conditions. This is the key architectural difference from all prior
systems: the formalisation boundary is not a gap in the encoding but a typed,
first-class object.

% ---------------------------------------------------------------
\section{Case Studies}
% ---------------------------------------------------------------

\subsection{IRC \S163(h): Mortgage Interest Deduction}

IRC~\S163(h) (26~U.S.C.~\S163(h)) is the provision of the US Internal Revenue
Code governing the deductibility of qualified residence interest. It was chosen
because Lawsky~\cite{lawsky2017logic} used it as the gold standard case study
for default logic applied to statutes. The rule structure consists of three
rules, four defeat relations forming a three-level chain, and three open-texture
terms.

The three rules in \texttt{LegalLean.CaseStudies.IRC163h} are:
\begin{itemize}
  \item \textbf{R1} (\texttt{R1\_BaseDeduction}): Taxpayers are permitted to
    deduct interest on acquisition indebtedness secured by a qualified
    residence. Priority level~1 (``statute'').
  \item \textbf{R2} (\texttt{R2\_TCJAProhibition}): Home equity interest
    deduction is prohibited for taxable years beginning after 15~December 2017
    (the TCJA). Priority level~2 (``statute (TCJA 2017, \textit{lex posterior})'').
  \item \textbf{R3} (\texttt{R3\_GrandfatherRule}): Acquisition indebtedness
    outstanding on 13~October 1987 is grandfathered from current dollar limits.
    Priority level~3 (``statute (\textit{lex specialis}, grandfather clause)'').
\end{itemize}

The four defeat relations are: R2 defeats R1 (\textit{lex posterior}); E1.1
(dollar limit) defeats R1; E1.2 (home equity dollar limit) defeats R1; and R3
defeats E1.1. The third of these creates a three-level defeasibility chain:
R3 defeats E1.1, which defeats R1 -- defeat of the limit restores the base
permission. This chain cannot be represented by a flat priority ordering; it
requires the full defeasibility machinery.

Open-texture terms are encoded as axioms with \texttt{Vague} instances.
\texttt{IsQualifiedResidence} is declared as an axiom, with the instance
documenting the core (principal residence plus one other property used as a
residence) and the penumbra (mobile homes, timeshares, boats with habitability
facilities, fractional ownership). \texttt{IsSubstantialImprovement} receives
the same treatment. \texttt{MixedUseAllocation} receives a
\texttt{Discretionary} instance referencing IRS tracing rules under
Treas.~Reg.~\S1.163-8T.

The compliance function \texttt{isDeductible} returns a
\texttt{FormalisationBoundary Prop}. For unlicensed entities, post-TCJA home
equity, and unsecured debt, it returns \texttt{formal~False} -- machine-verified
rejection. For acquisition indebtedness within limits, it returns
\texttt{boundary} with the authority ``IRS / Tax Court''.

\subsubsection{The Monotonicity Break Theorem.}

The headline result of this case study is \texttt{monotonicity\_break\_general}:

\begin{lstlisting}[language=lean4]
theorem monotonicity_break_general
    (payment : InterestPayment) (status : FilingStatus)
    (h_individual : payment.paidByIndividual = true)
    (h_secured    : payment.indebtedness.securedByResidence = true)
    (h_purpose    : payment.indebtedness.purpose = IndebtednessPurpose.other)
    (h_within     : withinHomeEquityLimit payment.indebtedness status = true)
    : (∃ r a, isDeductible payment ⟨2017⟩ status =
                FormalisationBoundary.boundary r a) ∧
      (isDeductible payment ⟨2018⟩ status =
                FormalisationBoundary.formal False) := by
  constructor
  · simp [isDeductible, h_individual, h_secured, h_purpose, h_within]
  · simp [isDeductible, h_individual, h_secured, h_purpose]
\end{lstlisting}

This theorem states, for all home equity interest payments by individual
taxpayers on secured property within the applicable dollar limits: the pre-TCJA
outcome (year~=~2017) is \texttt{boundary} (potentially deductible; human
determination of ``qualified residence'' required), while the post-TCJA outcome
(year~=~2018) is \texttt{formal~False} (automatically rejected;
machine-verified). The TCJA added a rule to the legal system that removed a
previously available permission. Lean~4 verifies this as a structural property
of the encoding, not merely a test case.

A companion corollary \texttt{pre\_tcja\_not\_auto\_rejected} confirms the
strictness of the break: the pre-TCJA outcome is provably not
\texttt{formal~False}, making the transition genuinely from ``possibly
permitted'' to ``certainly prohibited.''

It is worth being explicit about what this theorem proves and does not prove. It
proves a structural property of the formalisation: that the \texttt{isDeductible}
function, as encoded, produces these two outcomes for these two inputs. It does
not prove that the formalisation is a complete representation of IRC~\S163(h) --
such a claim would require independent legal audit. The theorem's value is that
it is machine-checked: the outcome cannot diverge silently from the encoding.

The proof's brevity (two \texttt{simp} calls) does not indicate a trivially
encoded result. The \texttt{simp} tactic unfolds \texttt{isDeductible},
substitutes the hypotheses, and evaluates the \texttt{if} branches. The
non-triviality is in the \textit{universality}: the theorem quantifies over all
\texttt{InterestPayment} values and all \texttt{FilingStatus} values satisfying
the four preconditions. A test case checks one input; this theorem checks the
entire input space. The preconditions are exactly the conditions under which the
monotonicity break is legally relevant (individual taxpayer, secured home equity
debt within limits), and the universally quantified conclusion is that the legal
outcome changes at the 2018 boundary for every such payment. This structural
coverage is what distinguishes a verified theorem from a unit test, even when
both share the same tactic.

\subsection{O-1A Visa: ``3 of 8'' as Dependent Types}

The O-1A extraordinary ability visa (8~CFR~\S204.5(h)) requires demonstrating
``extraordinary ability'' by meeting at least three of eight evidentiary
criteria. This is a natural dependent-type problem: the ``3 of 8'' threshold is
a cardinality constraint on a set of satisfied criteria, and Lean~4 can encode
this constraint at the type level.

The eight criteria are defined as an inductive type \texttt{Criterion} with
\texttt{DecidableEq}, allowing duplicate-free lists. Since the case study avoids
Mathlib, cardinality is encoded without \texttt{Finset}:

\begin{lstlisting}[language=lean4]
inductive AllDistinct {α : Type} [DecidableEq α] : List α → Prop
  | nil  : AllDistinct []
  | cons : ∀ {x : α} {xs : List α}, x ∉ xs → AllDistinct xs
               → AllDistinct (x :: xs)

structure MeetsThreeCriteria (a : Applicant) where
  satisfied : List Criterion
  distinct  : AllDistinct satisfied
  cardinal  : satisfied.length ≥ 3
  evidence  : ∀ c ∈ satisfied, LegalLean.EvidenceFor (criterionProp a c)
\end{lstlisting}

A value of type \texttt{MeetsThreeCriteria a} can only be constructed if one
provides a duplicate-free list of at least three criteria together with
\texttt{EvidenceFor} witnesses for each. The type checker enforces the threshold
at compile time. This is precisely the propositions-as-types reading of ``the
applicant meets the threshold'': inhabiting the type is proving the claim.

The Kazarian two-part test (\textit{Kazarian v.\ USCIS}, 596~F.3d~1115,
9th~Cir.~2010) is encoded by \texttt{determineEligibility}, which returns
\texttt{FormalisationBoundary.mixed}: the \texttt{formalPart} records that Lean
verified the criteria count; the \texttt{humanPart} names the USCIS adjudicator
and the final merits determination as the non-formalisable component. The
theorem \texttt{eligibility\_always\_mixed} proves that
\texttt{determineEligibility} always returns \texttt{mixed} -- never
\texttt{formal} or \texttt{boundary} alone -- which is a structural consequence
of the Kazarian framework: Part~1 is always machine-checkable; Part~2 is always
not.

\subsection{TCP Code: Deontic Spectrum and DDL Comparison}

The third case study encodes a kernel of the Australian Telecommunications
Consumer Protections Code (C628:2012), chosen for direct comparison with
Governatori's DDL encoding of the same regulatory
text~\cite{governatori2005representing,governatori2019compliance}. Three
provisions are encoded: TCP~3.2 (contact details obligation), TCP~5.1 (Critical
Information Summary obligation), and TCP~8.2 (excessive early termination fee
prohibition). Each has a \textit{lex specialis} defeat: small provider
exemption, renewal exemption, and device repayment carve-out respectively.

The three provisions demonstrate the full deontic spectrum: TCP~3.2 imposes an
obligation (providers must publish contact details); TCP~8.2 imposes a
prohibition (providers must not charge excessive ETFs); and the defeat rules
introduce permissions (small providers may use simplified requirements; device
repayment fees are allowed). The theorems
\texttt{contact\_obligation\_and\_exemption\_are\_deontic\_distinct} and
\texttt{etf\_prohibition\_and\_carveout\_are\_deontic\_distinct} prove that the
base rule and its defeating exception always have distinct modalities.

The \texttt{exemption\_rules\_outrank\_base\_rules} theorem verifies that
defeating rules have strictly higher priority levels than the rules they defeat
across all three provisions:

\begin{lstlisting}[language=lean4]
theorem exemption_rules_outrank_base_rules :
    R1_SmallProviderExemption.priority.level > R1_ContactObligation.priority.level ∧
    R2_RenewalExemption.priority.level > R2_CriticalInfoSummary.priority.level ∧
    R3_DeviceRepaymentCarveout.priority.level > R3_ExcessiveETFProhibited.priority.level
\end{lstlisting}

This is a structural soundness property: a defeat relation where the defeating
rule had lower or equal priority would be incoherent.

The DDL encoding of the same provisions is documented as comments in the source
file. For example, the contact obligation in DDL syntax:

\begin{lstlisting}[style=plain]
d ⊗ (isLicensed(P)) ⇒ O(publishContactDetails(P))
r1 :- licensed(P), not small_provider(P).
r1_exc :- small_provider(P).
r1_exc > r1
\end{lstlisting}

The LegalLean encoding differs in four respects: (1)~\texttt{FormalisationBoundary}
distinguishes \texttt{prominently\_published} from \texttt{is\_licensed} as
requiring human determination, rather than treating both as ordinary literals;
(2)~\texttt{LegalRule~(\(\alpha\)~:~Type)} parameterises rules by the entity
type, producing a typed rather than flat signature; (3)~\texttt{EvidenceFor~P}
requires a proof witness rather than mere assertion; (4)~the \texttt{Vague} and
\texttt{Discretionary} instances carry auditable metadata that is
machine-readable, not just a comment.

% ---------------------------------------------------------------
\section{Evaluation}
% ---------------------------------------------------------------

\subsection{Theorem Count and Properties}

The system produces 44 machine-verified theorems with zero \texttt{sorry}
placeholders across the three case studies. It is important to be honest about
the nature of these theorems. Some are structural unit tests:
\texttt{non\_individual\_never\_deducts} is proved by
\texttt{simp~[isDeductible,~h]} -- a trivial path trace through a function.
These theorems are valuable as regression tests but do not constitute substantive
contributions in isolation.

Others are more substantive. \texttt{monotonicity\_break\_general} requires
reasoning about two distinct legal regimes and proving a structural divergence.
\texttt{defeat\_chain\_exists} proves that the defeat relation in IRC~\S163(h)
has genuine depth -- R3 defeats E1.1, which defeats R1, but R3 does not directly
defeat R1 -- capturing a legally meaningful subtlety about the grandfather
clause. \texttt{no\_mixed\_outcomes} proves by full case analysis that
IRC~\S163(h)'s \texttt{isDeductible} function never returns a \texttt{mixed}
result, establishing a binary structure as a property of the statute rather than
a design choice.

A faithfulness audit identified three areas where the encoding simplifies the
actual law: (1)~date comparisons use ISO~8601 string ordering rather than proper
date arithmetic; (2)~\texttt{criterionProp} now maps each O-1A criterion to its
own distinct axiom (e.g.\ \texttt{HasExtraordinaryAbility},
\texttt{SalaryIsHigh}, \texttt{ContributionsAreMajor}), with each carrying a
separate \texttt{Vague} or \texttt{Discretionary} instance that documents the
specific core and penumbra for that criterion; (3)~the IRC~\S163(h) encoding
does not handle the full debt tracing rules of Treas.~Reg.~\S1.163-8T, which
require temporal tracking of how loan proceeds were used.

\subsection{Comparison with Prior Systems}

Table~\ref{tab:comparison} compares LegalLean against four prior systems across
seven dimensions on a 1--5 scale. As noted below, these scores should be treated
with caution.

\begin{table}[t]
\caption{Comparison of LegalLean with prior legal formalisation systems.
Scores are on a 1--5 scale; higher is better. Scores reflect the dimensions on
which LegalLean makes a genuine claim, not an objective ranking.}
\label{tab:comparison}
\begin{tabular}{lccccc}
\hline
\textbf{Dimension} & \textbf{LegalLean} & \textbf{L4L} & \textbf{DDL} &
  \textbf{Catala} & \textbf{LegalRuleML} \\
\hline
Compositional semantics         & 5 & 3 & 4 & 4 & 2 \\
Deontic modalities              & 5 & 4 & 5 & 2 & 4 \\
Dependent types                 & 5 & 5 & 1 & 1 & 1 \\
Machine-checkable guarantees    & 5 & 3 & 3 & 4 & 1 \\
Open texture encoding           & 5 & 3 & 4 & 2 & 2 \\
Provenance tracking             & 5 & 5 & 4 & 3 & 4 \\
Temporal reasoning              & 5 & 3 & 4 & 3 & 4 \\
\hline
\textbf{Total}                  & \textbf{35} & \textbf{26} & \textbf{25} &
  \textbf{19} & \textbf{18} \\
\hline
\end{tabular}
\end{table}

These scores were generated as part of a structured comparison exercise and must
be interpreted with appropriate scepticism. The scores for LegalLean are perfect
(5/5 on every dimension), which reflects the dimensions being chosen partly to
highlight LegalLean's strengths. A reviewer who weighted practical adoption, tool
maturity, or non-expert usability would produce different rankings.

\textbf{DDL}~\cite{governatori2005representing,governatori2019compliance}: strong
on deontic modalities and temporal reasoning; the most mature industrial-strength
defeasible logic system for law. Its limitation is that it treats all conditions
as first-order literals -- \texttt{prominently\_published(P)} is no different
from \texttt{is\_licensed(P)} in the DDL signature. There is no mechanism to
mark the former as requiring institutional determination.

\textbf{Catala}~\cite{merigoux2021catala}: designed specifically for tax and
social benefit statutes, with excellent practical coverage. It computes legal
outcomes rather than proving properties about them: Catala programs execute, they
do not produce proof certificates. It has no representation of open texture. In
LegalLean, \texttt{isDeductible} returns
\texttt{FormalisationBoundary.boundary~"qualified~residence"~"IRS~/~Tax~Court"}
-- the system records that it cannot determine this condition and names the
authority responsible. The distinction is not that LegalLean ``handles''
vagueness better, but that it represents the inability to handle it as a typed,
auditable object rather than a silent gap.

\textbf{L4L}~\cite{hou2021l4l}: the closest prior work in spirit. L4L uses
dependent types and produces proofs, but SMT-backed verification means the proof
certificates are not auditable proof terms in the Curry-Howard sense. Its
treatment of open texture relies on the SMT solver's inability to discharge
certain goals rather than a typed annotation.

\textbf{LegalRuleML}~\cite{palmirani2018modelling}: a rich XML ontology with
sophisticated modelling of deontic modalities and temporal applicability. It is a
knowledge representation language, not a reasoning system -- it provides no
automated inference or proof production.

\subsection{Limitations}

\textbf{Scalability}: the three case studies cover nine provisions and fewer than
300 lines of Lean code each. Scaling to a complete statute would require
substantially more proof engineering.

\textbf{No LLM-to-Lean pipeline}: the framework does not yet include an
automated formalisation pipeline. The transducer half of the architecture --
converting statutory text to Lean~4 -- is currently manual.

\textbf{Axiom dependence}: the open-texture axioms are genuine axioms that
cannot be proved by Lean's kernel. The framework is only as sound as the legal
judgements behind the \texttt{Vague} and \texttt{Discretionary} instances.

\textbf{No temporal algebra}: the TCJA cutoff is encoded as a numeric comparison
(\texttt{year.year >= 2018}), not via a proper temporal type with duration,
retroactivity, and sunset provisions. A complete temporal logic layer remains
future work.

\textbf{Sorry-bypass gap}: the type system does not prevent the construction
\texttt{FormalisationBoundary.formal~sorry}. A careless user could mark a
proposition as formally verified using Lean's \texttt{sorry} axiom. The
VERIFICATION.md protocol addresses this by requiring
\texttt{grep~-r~sorry} as part of the verification layer, but the type system
itself does not enforce it. This is an engineering limitation, not a theoretical
one.

\textbf{Normative risk of formalisation}: by encoding Hart's open texture as a
typed object, this work makes the formalisation boundary tractable and auditable.
This carries a normative risk that Lessig's ``code is law'' tradition identifies:
making the boundary typed may create a false sense of completeness. A user who
sees \texttt{FormalisationBoundary.boundary~"requires~IRS~determination"~"IRS~/~Tax~Court"}
might treat the annotation as sufficient engagement with the underlying vagueness,
when in fact the core/penumbra analysis in the \texttt{Vague} instance is one
interpretation among several. The system does not resolve open texture; it
classifies and annotates it. Users must understand that the annotation is a map
of the boundary, not a resolution of it. This limitation is intrinsic to any
formalisation of vagueness and should be communicated clearly in any deployment
context.

% ---------------------------------------------------------------
\section{Related Work}
% ---------------------------------------------------------------

\textbf{Defeasible Deontic Logic
(DDL)}~\cite{governatori2005representing,governatori2019compliance}: the most
complete prior work on defeasible reasoning in law. It provides a proof theory
for defeasible rules with deontic modalities and has been applied to compliance
checking in eHealth, business contracts, and regulatory codes. LegalLean's
advance over DDL is not in defeasibility machinery but in (1)~machine-checkable
proof certificates, (2)~the \texttt{FormalisationBoundary} typing of open
texture, and (3)~dependent types enabling evidence-as-inhabitants.

\textbf{Catala}~\cite{merigoux2021catala}: the most practical prior system for
statutory formalisation, with a growing library of French, US, and Polish tax
law. Its design priority is executability: a Catala specification is a runnable
programme that computes legal outcomes from facts. LegalLean's priority is
verifiability. The two approaches are complementary rather than competitive. A
concrete example illustrates this complementarity: in IRC~\S163(h), the
condition ``property is a qualified residence'' is vague in Hart's sense (the
penumbra includes mobile homes, timeshares, and boats). In Catala, this
condition must either be encoded as a Boolean input (the user asserts it) or
omitted (the programme raises an error for unhandled cases). In LegalLean,
\texttt{isDeductible} returns \texttt{FormalisationBoundary.boundary~"qualified~residence"~"IRS~/~Tax~Court"}
-- the system records that it cannot determine this condition, names the
authority responsible, and the \texttt{Vague} instance on
\texttt{IsQualifiedResidence} documents the core and penumbra for audit. The
distinction is not that LegalLean ``handles'' vagueness better, but that it
represents the inability to handle it as a typed, auditable object rather than a
silent gap or a runtime error.

\textbf{L4L}~\cite{hou2021l4l}: the work closest in motivation to LegalLean.
The primary technical difference is that L4L's verification layer relies on SMT,
while LegalLean's relies on Lean~4's kernel verifier. SMT provides more
automation; Lean~4's kernel provides more auditability. L4L does not include a
typed treatment of open texture.

\textbf{LegalRuleML}~\cite{palmirani2018modelling}: extends RuleML with deontic
operators, temporal metadata, and defeasibility annotations. It is the most
expressive knowledge representation format for legal rules but provides no
automated reasoning. LegalLean and LegalRuleML address different problems:
LegalLean verifies reasoning; LegalRuleML represents knowledge for interchange.

\textbf{Hohfeldian frameworks}: Prior work by
McCarty~\cite{mccarty1977reflections} and
Boella~et~al.~\cite{boella2006introduction} applies Hohfeld's fundamental legal
relations in formal settings. LegalLean's three-modality deontic type is
intentionally minimal; a Hohfeldian extension is straightforward but not required
for the case studies.

\textbf{Logic programming approaches}: Kowalski and Sergot's Event
Calculus~\cite{shanahan1999event}, together with Sergot~et~al.'s British
Nationality Act encoding~\cite{sergot1986british}, established the research
programme of encoding law in logic programming. That tradition addressed
defeasibility before DDL but did not handle open texture or produce
machine-checked proofs.

% ---------------------------------------------------------------
\section{Conclusion and Future Work}
% ---------------------------------------------------------------

LegalLean demonstrates that Lean~4 dependent types provide a viable verification
layer for legal reasoning systems. The key result is not the theorem count but
the architecture: a system in which Hart's open texture is a first-class typed
object rather than an external limitation, and in which the distinction between
``machine-verified'' and ``requires human determination'' is enforced by the type
checker rather than stated in documentation.

The \texttt{monotonicity\_break\_general} theorem illustrates what
machine-checked legal reasoning makes possible: a formal, auditable certificate
that the TCJA 2017 removed a deduction permission for all home equity interest
payments, not merely for specific test cases. This kind of structural certificate
-- not an answer, but a proof about the legal system itself -- is what
distinguishes a verifier from a transducer.

Several directions remain open. First, the transducer layer: language models are
the natural tool for converting statutory text to Lean~4 stubs, with the type
checker providing feedback to guide the formalisation. Second, zero-knowledge
proofs for compliance: a party wishing to attest compliance without disclosing
underlying facts could, in principle, use ZK proofs over the formal structure.
Third, temporal algebra: a proper treatment of statutory effective dates, sunset
provisions, and retroactivity requires a temporal type extending the current
\texttt{TaxYear} structure. Fourth, scaling studies: encoding a complete statute
rather than a selected kernel would test whether the proof engineering effort is
tractable in practice.

The verification protocol that emerges from this work has three layers: the type
checker verifies structural properties; the \texttt{FormalisationBoundary}
annotations identify where institutional determination is required; and legal
practitioners review the open-texture instances to ensure the core/penumbra
analysis is sound. Only the first layer is automated. The claim is not that legal
reasoning can be automated but that the parts that can be verified should be, and
the parts that cannot should be typed rather than silently omitted.

\begin{credits}
\subsubsection{\discintname}
The author has no competing interests to declare that are relevant to the content
of this article.
\end{credits}

% ---- Bibliography ----
\bibliographystyle{splncs04}
\bibliography{references}

\end{document}
